% ============================================================
%  LaTeX Beamer 演示模板 —— 完整功能教学版
%  编译引擎：XeLaTeX（支持中文，UTF-8 编码）
%  编译命令：xelatex beamer-template.tex
%  屏幕比例：16:9（宽屏）
%  作者：GitHub Copilot
%  日期：2026-05-18
%
%  本模板涵盖 Beamer 的绝大多数常用功能，每个部分均有
%  详细的注释说明，适合 LaTeX 初学者系统学习 Beamer。
% ============================================================

% ============================================================
%  第 1 部分：文档类（Document Class）
% ============================================================
%
%  Beamer 是 LaTeX 的演示文稿文档类。
%  关键选项说明：
%   - UTF8          : 启用 UTF-8 编码支持（配合 xelatex）
%   - 11pt          : 正文默认字号，可选 8pt/9pt/10pt/11pt/12pt/14pt/17pt/20pt
%   - aspectratio=169 : 屏幕比例，可选 43(4:3) / 169(16:9) / 1610(16:10) / 32(3:2) / 149(14:9)
%   - handout       : （已注释）生成讲义版本，无动画叠层
%   - t / trans     : （可选）使用 transparency 模式
%   - compress      : （可选）压缩导航栏
%
\documentclass[UTF8,11pt,aspectratio=169]{beamer}
% 如需生成讲义版本（去掉逐层动画，所有叠层合并），取消下一行注释：
% \documentclass[UTF8,11pt,aspectratio=169,handout]{beamer}

% ============================================================
%  第 2 部分：中文与字体设置
% ============================================================
%
%  xeCJK 是在 XeLaTeX 下处理中日韩文字的标准方案。
%  字体名称是 macOS 系统自带的——Windows/Linux 用户需
%  替换为对应平台的字体名。
%
\usepackage{xeCJK}                          % 中日韩文字支持
\setCJKmainfont{Songti SC}                  % 正文中文字体：宋体（macOS）
% \setCJKmainfont{SimSun}                   % Windows：宋体
\setCJKsansfont{Heiti SC}                   % 无衬线中文字体：黑体（macOS）
% \setCJKsansfont{SimHei}                   % Windows：黑体
\setCJKmonofont{STFangsong}                 % 等宽中文字体：仿宋（macOS）
% \setCJKmonofont{FangSong}                 % Windows：仿宋

% ============================================================
%  第 3 部分：主题（Theme）与配色（Color Theme）
% ============================================================
%
%  Beamer 通过三层设定来控制外观：
%   1. 整体主题（Theme）      —— 布局、导航栏、标题样式
%   2. 配色主题（Color Theme）—— 各元素的颜色
%   3. 字体主题（Font Theme） —— 字体组合
%
%  ---- 常用整体主题（取消注释即可切换） ----
% \usetheme{Madrid}            % 经典蓝色导航栏，简洁专业（当前使用）
% \usetheme{Berlin}          % 顶部有导航点的横向条
% \usetheme{Copenhagen}      % 类似 Madrid 但更柔和
% \usetheme{Singapore}       % 极简风格
% \usetheme{Warsaw}          % 有阴影效果的大导航栏
% \usetheme{AnnArbor}        % 蓝色与黄色搭配
% \usetheme{Pittsburgh}      % 右侧竖条
% \usetheme{Rochester}       % 宽横幅导航
% \usetheme{Boadilla}        % 简洁优雅
\usetheme{Darmstadt}       % 带有平滑导航栏
%
%  特别推荐：metropolis 是一个现代扁平化主题（需单独安装）
% \usetheme{metropolis}      % 需先安装：https://github.com/matze/mtheme

% ---- 配色主题（与 Madrid 配合良好） ----
\usecolortheme{whale}        % 深蓝色调（当前使用）
% 其他常用配色：dolphin / seahorse / orchid / crane / beaver

% ---- 字体主题 ----
\usefonttheme{professionalfonts}  % 使用系统安装的字体，而非 Beamer 默认替换
% 其他可选：serif / structurebold / structureitalicserif

% ---- 自定义颜色（可用于文本、block、表格等） ----
\definecolor{darkgreen}{RGB}{0,128,0}       % 深绿色
\definecolor{darkred}{RGB}{180,0,0}         % 深红色
\definecolor{darkblue}{RGB}{0,0,139}        % 深蓝色
\definecolor{gold}{RGB}{200,150,0}          % 金色
\definecolor{mygray}{RGB}{245,245,245}      % 浅灰色（代码块背景）

% ---- 图标设置（调整导航符号显示） ----
% \setbeamertemplate{navigation symbols}{}    % 隐藏右下角导航符号（更简洁）
% \setbeamertemplate{navigation symbols}[vertical]  % 垂直排列
\setbeamertemplate{navigation symbols}[horizontal] % 水平排列（默认）

% ---- 调整标题页和帧标题的样式 ----
% \setbeamercolor{frametitle}{fg=white,bg=darkblue}  % 帧标题白字蓝底

% ============================================================
%  第 4 部分：常用宏包
% ============================================================
%
%  以下按功能分组列出，注释说明了每个宏包的核心用途。
%

% ---------- 图形与绘图 ----------
\usepackage{graphicx}        % 插入外部图片 \includegraphics
\usepackage{tikz}            % 矢量绘图引擎
\usetikzlibrary{
    positioning,             % 相对定位（above=of, right=of 等）
    calc,                    % 坐标计算（$...$）
    shapes.geometric,        % 几何形状节点（椭圆、菱形等）
    arrows.meta,             % 新版箭头样式（Stealth, Latex 等）
    backgrounds,             % 背景层（用于衬底）
    fit,                     % 包围盒（fit 多个节点）
    chains                   % 链式节点排列
}
% \usepackage{subcaption}      % ⚠️ 不要在 Beamer 中使用 subcaption！与 Beamer 冲突

% ---------- 表格 ----------
\usepackage{booktabs}        % 三线表：\toprule \midrule \bottomrule
\usepackage{tabularx}        % 可伸缩列宽的表格
\usepackage{multirow}        % 纵向合并单元格
\usepackage{array}           % 增强列格式（>{\centering}m{2cm} 等）
\usepackage[table]{xcolor}   % 表格行/列/单元格着色

% ---------- 数学 ----------
\usepackage{amsmath,amssymb} % AMS 数学宏包（公式对齐、数学符号）
\usepackage{mathtools}       % 增强数学排版（\coloneqq, \DeclarePairedDelimiter 等）

% ---------- 代码与列表 ----------
\usepackage{listings}        % 源代码语法高亮
\lstset{                     % 全局代码块样式设定
    basicstyle=\ttfamily\small,          % 代码字体：等宽小号
    keywordstyle=\color{blue}\bfseries,  % 关键字：蓝色加粗
    commentstyle=\color{darkgreen},      % 注释：深绿色
    stringstyle=\color{darkred},         % 字符串：深红色
    numbers=left,                        % 行号在左侧
    numberstyle=\tiny\color{gray},       % 行号样式
    backgroundcolor=\color{mygray},      % 背景色
    breaklines=true,                     % 长行自动换行
    frame=single,                        % 单线边框
    tabsize=2                            % Tab 宽度
}
% \usepackage{enumitem}        % ⚠️ 不要在 Beamer 中加载 enumitem！会与 Beamer 的列表机制冲突

% ---------- 链接与引用 ----------
\usepackage{hyperref}        % 超链接支持
\hypersetup{
    colorlinks=true,         % 链接着色（而非方框）
    linkcolor=darkblue,      % 内部链接颜色
    urlcolor=darkblue        % URL 链接颜色
}

% ---------- 其他工具 ----------
\usepackage{multirow}        % 已在表格部分；确保加载
\usepackage{bookmark}        % PDF 书签（替代 hyperref 的书签功能）
\usepackage{twemojis}

% ============================================================
%  第 5 部分：文档信息（标题页数据）
% ============================================================
%
%  这些命令定义的信息会出现在 \titlepage 中。
%  方括号 [] 中的是短版本，用于页脚或导航栏。
%
\title[Beamer 模板教学]{LaTeX Beamer 完整功能演示模板}
%      [短标题（可选）] {完整标题（必需）}

\subtitle{从入门到精通 —— 涵盖幻灯片制作的所有核心功能}
% 副标题（可选），会以小号字显示在标题下方

\author[GitHub Copilot]{GitHub Copilot}
%      [短作者名]        {完整作者名}

\date{\today}                                % 日期：\today 表示编译当天
% \date{2026年5月18日}                       % 也可手动指定

\institute{数字IC 设计}                      % 机构/单位（可选）
% \institute[短机构名]{完整机构名}

\titlegraphic{\vspace*{-1cm}}                % 标题页底部图形（留空以去掉默认线）
% \titlegraphic{\includegraphics[width=2cm]{logo.png}}  % 也可插入 Logo

% ============================================================
%  第 6 部分：自定义命令与环境
% ============================================================
%
%  自定义命令可以大幅提高写作效率。
%

% ---- 6.1 自定义快捷命令 ----
\newcommand{\beamer}{Beamer}                 % 用 \beamer 替代 Beamer（保持格式）
\newcommand{\latex}{\LaTeX}                  % 简写
\newcommand{\xelatex}{Xe\LaTeX}              % 简写
\newcommand{\important}[1]{\textcolor{red}{\textbf{#1}}}
                                             % 以红色加粗突出显示文字

% ---- 6.2 自定义 Block 环境 ----
%  创建语义化的彩色框，在文档中用 \begin{keypoint}...\end{keypoint} 调用
\newenvironment{keypoint}{
    \begin{block}{\textbf{\twemoji{key} 关键结论}}
}{
    \end{block}
}

\newenvironment{warningblock}{
    \begin{alertblock}{\textbf{\twemoji{warning}️ 注意}}
}{
    \end{alertblock}
}

\newenvironment{tip}{
    \begin{exampleblock}{\textbf{\twemoji{light bulb} 提示}}
}{
    \end{exampleblock}
}

% ---- 6.3 自定义 Beamer 模板（进阶） ----
%  修改脚注显示格式
\setbeamertemplate{footline}[frame number]
% 其他可选：infolines / split / text line

% ---- 6.4 表格中数字按小数点对齐的命令 ----
%  来自 siunitx 宏包（如需要可取消注释）
% \usepackage{siunitx}
% \newcommand{\numtab}[1]{\num{#1}}

% ============================================================
%  第 7 部分：目录（Table of Contents）自动化
% ============================================================
%
%  以下设置让 Beamer 在每个 section 开始前自动显示目录页，
%  并高亮当前章节（currentsection）。
%
\AtBeginSection[]{
    \begin{frame}{目录}
        %  \tableofcontents 自动生成目录
        %  currentsection 选项：只显示当前 section 并高亮
        \tableofcontents[currentsection]
    \end{frame}
}

%  如果也想在 subsection 前显示目录：
% \AtBeginSubsection[]{
%     \begin{frame}{目录}
%         \tableofcontents[currentsection,currentsubsection]
%     \end{frame}
% }

% ============================================================
%  文档正文开始
% ============================================================
\begin{document}

% ============================================================
%  幻灯片 0：标题页
% ============================================================
\begin{frame}
    \titlepage
    %  \titlepage 自动使用 \title, \subtitle, \author, \date, \institute
\end{frame}

% ============================================================
%  幻灯片 1：概览目录页
% ============================================================
\begin{frame}{目录}{本演示的结构概览}
    % 第一个 {} 是帧标题，第二个 {} 是帧副标题（可选）
    \tableofcontents
    %  不带任何选项：显示完整目录
\end{frame}

% ============================================================
%  SECTION 1：基础幻灯片结构
% ============================================================
\section{基础幻灯片结构}

% ---- 1.1 文本与段落 ----
%  [fragile] 必需：因为 frame 中使用了 \verb 命令
\begin{frame}[fragile]{文本与段落}{基础排版要素}
    %
    %  \frametitle{...} 和 \framesubtitle{...} 是 frame 环境中
    %  替代 {} 传参的另一种写法。
    %

    % --- 普通文本 ---
    这是普通正文文本。Beamer 默认使用无衬线（Sans-serif）字体，
    适合屏幕投影阅读。

    % --- 段落间距 ---
    这是第二段。
    在 Beamer 中，空行表示新段落，但段落间距通常比 article 类小。

    % --- 强调文本 ---
    \begin{itemize}
        \item \textbf{粗体}：使用 \verb|\textbf{...}|
        \item \textit{斜体}：使用 \verb|\textit{...}|
        \item \texttt{等宽}：使用 \verb|\texttt{...}|
        \item \textcolor{blue}{颜色}：使用 \verb|\textcolor{blue}{...}|
        \item \alert{警告色}：使用 \verb|\alert{...}|（自动使用主题警示色）
        \item \important{自定义高亮}：使用我们定义的 \verb|\important{...}|
    \end{itemize}
\end{frame}

% ---- 1.2 列表 ----
\begin{frame}{列表环境}{有序列表与无序列表}
    %
    %  Beamer 中列表是逐步显示的（默认 pause 逐条）。
    %  如需一次性显示全部，可用 [<+->] 或去掉逐层标记。
    %

    % --- 无序列表 (itemize) ---
    无序列表（itemize）：
    \begin{itemize}
        \item 第一项：itemize 默认使用三角形符号
        \item 第二项：嵌套列表
        \begin{itemize}
            \item 子项 1 —— 二级列表用短横线
            \item 子项 2
        \end{itemize}
        \item 第三项：可以包含公式 $E = mc^2$
    \end{itemize}

    % --- 有序列表 (enumerate) ---
    有序列表（enumerate）：
    \begin{enumerate}
        \item 步骤一：定义问题
        \item 步骤二：设计方案
        \begin{enumerate}
            \item 方案 A：简单直接
            \item 方案 B：复杂但灵活
        \end{enumerate}
        \item 步骤三：实施并验证
    \end{enumerate}
\end{frame}

% ---- 1.3 分栏布局 ----
\begin{frame}{分栏布局}{columns 环境的使用}
    %
    %  columns 是 Beamer 中最常用的布局工具。
    %  每一列的宽度用 \textwidth 的比例来指定。
    %

    \begin{columns}[T]  % [T] = 顶部对齐, [c] = 居中, [b] = 底部对齐, [onlytextwidth]
        %
        % ---- 左栏 50% ----
        \begin{column}{0.5\textwidth}
            \textbf{左栏内容}：
            \begin{itemize}
                \item 分栏是并排展示信息的
                \item 最佳方式
                \item 常用于图文并排
                \item 或中英文对照
                \item 或方案对比
            \end{itemize}
        \end{column}
        %
        % ---- 右栏 50% ----
        \begin{column}{0.5\textwidth}
            \textbf{右栏内容}：
            \begin{enumerate}
                \item 左侧是列表
                \item 右侧也可以是列表
                \item 两侧内容完全独立
                \item 排版灵活度极高
            \end{enumerate}
            \vspace{0.3cm}
            \centering
            \alert{中间也可以放强调文字}
        \end{column}
    \end{columns}

    \vspace{0.5cm}
    % 分栏下方仍然可以放内容
    \hrule
    \vspace{0.3cm}
    \small 提示：列宽之和不超过页面文本宽度，否则会溢出。
\end{frame}

% ---- 1.4 三栏布局 ----
\begin{frame}{三栏布局}{灵活的多栏设计}
    \begin{columns}[T]
        \begin{column}{0.3\textwidth}
            \centering
            \textbf{列 1} \\
            \vspace{0.3cm}
            宽度 30\%，适合放图片或简短要点。
        \end{column}
        \begin{column}{0.4\textwidth}
            \centering
            \textbf{列 2（最宽）} \\
            \vspace{0.3cm}
            宽度 40\%，适合放主要内容。
            可以包含较长的文本描述。
        \end{column}
        \begin{column}{0.3\textwidth}
            \centering
            \textbf{列 3} \\
            \vspace{0.3cm}
            宽度 30\%，数据或补充说明。
        \end{column}
    \end{columns}
\end{frame}

% ============================================================
%  SECTION 2：块环境
% ============================================================
\section{块环境（Blocks）}

% ---- 2.1 三种标准 Block ----
\begin{frame}{三种标准 Block}{block / alertblock / exampleblock}
    %
    %  Block 是 Beamer 的核心视觉元素，用于将内容组织在
    %  带标题的彩色框内。
    %

    % --- 普通 Block（蓝色） ---
    \begin{block}{普通 Block}
        这是普通 block 环境，标题背景使用主题主色调。
        适合展示一般性结论或定义。
    \end{block}

    % --- 警示 Block（红色） ---
    \begin{alertblock}{警示 Block}
        这是 alertblock 环境，标题背景使用警示色（红色）。
        适合展示警告、注意事项。
    \end{alertblock}

    % --- 示例 Block（绿色） ---
    \begin{exampleblock}{示例 Block}
        这是 exampleblock 环境，标题背景使用示例色（绿色）。
        适合展示示例、正面案例。
    \end{exampleblock}

    % --- 自定义语义 Block（使用前文定义的 keypoint 环境） ---
    \begin{keypoint}
        使用自定义 block 环境可以实现语义化标记，
        比如这个"关键结论"块。
    \end{keypoint}
\end{frame}

% ---- 2.2 Block 内嵌复杂内容 ----
\begin{frame}{Block 内嵌复杂内容}{列表、公式、表格皆可放入 Block}
    %
    \begin{block}{Block 可包含多种内容}
        \begin{itemize}
            \item Block 内部可以放列表
            \item 也可以放数学公式：
            \[
                \int_{-\infty}^{\infty} e^{-x^2}\,dx = \sqrt{\pi}
            \]
            \item 甚至放小型表格：
            \begin{center}
            \begin{tabular}{c|c}
                \toprule
                参数 & 值 \\
                \midrule
                $\alpha$ & 0.05 \\
                $\beta$  & 0.95 \\
                \bottomrule
            \end{tabular}
            \end{center}
        \end{itemize}
    \end{block}
\end{frame}

% ============================================================
%  SECTION 3：表格
% ============================================================
\section{表格}

% ---- 3.1 三线表 ----
\begin{frame}{三线表}{booktabs 专业表格}
    %
    %  booktabs 提供 \toprule（顶线，粗）、\midrule（中线，细）、
    %  \bottomrule（底线，粗），比默认 \hline 更专业。
    %
    \begin{table}
        \centering
        \caption{各种 LaTeX 宏包的功能分类}
        \begin{tabular}{lll}
            % l = 左对齐, c = 居中, r = 右对齐, p{2cm} = 固定宽度
            \toprule
            \textbf{分类}    & \textbf{宏包}      & \textbf{功能}       \\
            \midrule
            表格             & booktabs           & 三线表专业排版      \\
                             & tabularx           & 可伸缩列宽          \\
                             & multirow           & 纵向合并单元格      \\
            \addlinespace   % 在行间增加额外间距
            图形             & graphicx           & 插入图片            \\
                             & tikz               & 矢量绘图            \\
                             & (minipage)         & 子图并排（替代 subcaption）\\
            \midrule
            数学             & amsmath            & 高级数学环境        \\
                             & amssymb            & 数学符号扩展        \\
            \bottomrule
        \end{tabular}
    \end{table}
\end{frame}

% ---- 3.2 合并单元格与着色 ----
\begin{frame}{合并单元格与着色}{multirow + 颜色}
    %
    \begin{table}
        \centering
        \caption{多行合并与行着色示例}
        \begin{tabular}{c|c|c}
            \toprule
            \textbf{阶段} & \textbf{任务}     & \textbf{负责人} \\
            \midrule
            \multirow{3}{*}{阶段一}
                & 需求分析  & 张三 \\
                & 方案设计  & 张三 \\
                & 评审      & 李四 \\
            \midrule
            \rowcolor{blue!10}  % 整行浅蓝背景
            \multirow{2}{*}{阶段二}
                & 开发      & 王五 \\
            \rowcolor{blue!10}
                & 测试      & 赵六 \\
            \bottomrule
        \end{tabular}
    \end{table}
\end{frame}

% ============================================================
%  SECTION 4：图片与图形
% ============================================================
\section{图片与图形}

% ---- 4.1 插入图片 ----
%  [fragile] 必需：因为 frame 中使用了 \verb 命令
\begin{frame}[fragile]{插入图片}{\texttt{\textbackslash includegraphics} 命令}
    %
    \begin{columns}[T]
        \begin{column}{0.5\textwidth}
            \begin{itemize}
                \item 支持的格式：PDF、PNG、JPG、EPS
                \item \verb|width|：指定宽度
                \item \verb|height|：指定高度
                \item \verb|scale|：缩放因子
                \item \verb|angle|：旋转角度
                \item \verb|trim|：裁剪
                \item \verb|clip|：裁剪生效
                \item \verb|keepaspectratio|：保持宽高比
            \end{itemize}
        \end{column}
        \begin{column}{0.5\textwidth}
            % 用一个彩色色块模拟图片（实际使用时替换为真实图片路径）
            \centering
            \begin{tikzpicture}
                \fill[blue!30] (0,0) rectangle (5,3);
                \node at (2.5,1.5) {\large\textbf{示例图片}};
                \node at (2.5,0.8) {\small (请替换为真实图片)};
            \end{tikzpicture}
            % 真实用法：
            % \includegraphics[width=0.9\textwidth,keepaspectratio]{figures/your-image.png}
        \end{column}
    \end{columns}
\end{frame}

% ---- 4.2 子图并排 ----
\begin{frame}{子图并排}{手动实现并排图片}
    %
    %  在 Beamer 中，figure 环境不会真正浮动，但提供 caption 功能。
    %  注意：Beamer 中不要使用 subcaption 宏包（与 Beamer 冲突）。
    %  手动用 minipage 或 columns 实现子图并排即可。
    %
    \begin{figure}
        \centering
        % --- 左侧子图（minipage） ---
        \begin{minipage}{0.45\textwidth}
            \centering
            \begin{tikzpicture}[scale=0.6]
                \fill[red!30]  (0,0) rectangle (4,3);
                \node at (2,1.5) {图 A};
            \end{tikzpicture}
            \caption*{(a) 左侧子图}
            %  \caption* 来自 Beamer，生成不带编号的标题
        \end{minipage}
        \hfill
        % --- 右侧子图（minipage） ---
        \begin{minipage}{0.45\textwidth}
            \centering
            \begin{tikzpicture}[scale=0.6]
                \fill[green!30] (0,0) rectangle (4,3);
                \node at (2,1.5) {图 B};
            \end{tikzpicture}
            \caption*{(b) 右侧子图}
        \end{minipage}
        \caption{两个子图并排展示（使用 minipage 实现）}
        \label{fig:sidebyside}
    \end{figure}
    %
    通过 \texttt{minipage} 或 \texttt{columns} 可以在 \texttt{figure} 内部
    并排放置子图，每个子图可以有自己的子标题。
\end{frame}

% ============================================================
%  SECTION 5：数学公式
% ============================================================
\section{数学公式}

% ---- 5.1 行内与行间公式 ----
%  [fragile] 必需：因为 frame 中使用了 \verb 命令
\begin{frame}[fragile]{行内与行间公式}{公式的两种模式}
    %
    % --- 行内公式（Inline Math） ---
    行内公式：$f(x) = x^2 + 2x + 1$ 嵌入在文本中。

    带编号的行间公式：
    \begin{equation}
        \label{eq:quadratic}
        f(x) = ax^2 + bx + c
    \end{equation}
    公式 \eqref{eq:quadratic} 是二次函数的标准形式。

    % --- 不带编号的行间公式 ---
    不带编号的行间公式（使用 \verb|\[ ... \]|）：
    \[
        \int_{0}^{\infty} \frac{\sin x}{x}\,dx = \frac{\pi}{2}
    \]
\end{frame}

% ---- 5.2 多行公式对齐 ----
\begin{frame}{多行公式对齐}{align 与 cases 环境}
    %
    %  align 环境：用 & 标记对齐点，\\ 换行
    %
    \begin{align}
        (a + b)^2 &= a^2 + 2ab + b^2             \label{eq:binomial1} \\
        (a - b)^2 &= a^2 - 2ab + b^2             \label{eq:binomial2} \\
        (a + b)(a - b) &= a^2 - b^2              \label{eq:binomial3}
    \end{align}

    %  cases 环境：分段函数
    分段函数示例：
    \[
        |x| =
        \begin{cases}
             x,  & x \geq 0 \\
            -x,  & x < 0
        \end{cases}
    \]
\end{frame}

% ---- 5.3 矩阵 ----
\begin{frame}{矩阵与线性代数}{pmatrix / bmatrix / vmatrix}
    %
    \begin{columns}[T]
        \begin{column}{0.45\textwidth}
            \textbf{pmatrix（圆括号）}：
            \[
            \begin{pmatrix}
                a_{11} & a_{12} \\
                a_{21} & a_{22}
            \end{pmatrix}
            \]
            \textbf{bmatrix（方括号）}：
            \[
            \begin{bmatrix}
                1 & 0 & 0 \\
                0 & 1 & 0 \\
                0 & 0 & 1
            \end{bmatrix}
            \]
        \end{column}
        \begin{column}{0.55\textwidth}
            \textbf{vmatrix（行列式）}：
            \[
            \begin{vmatrix}
                a & b \\
                c & d
            \end{vmatrix}
            = ad - bc
            \]
            带省略号的矩阵：
            \[
            A_{m \times n} =
            \begin{pmatrix}
                a_{11} & \cdots & a_{1n} \\
                \vdots & \ddots & \vdots \\
                a_{m1} & \cdots & a_{mn}
            \end{pmatrix}
            \]
        \end{column}
    \end{columns}
\end{frame}

% ============================================================
%  SECTION 6：代码展示
% ============================================================
\section{代码展示}

% ---- 6.1 代码块 ----
%  [fragile] 必需：因为 frame 中使用了 lstlisting 环境
\begin{frame}[fragile]{源代码展示}{listings 宏包的语法高亮}
    %
    \begin{block}{Python 代码示例}
        %  lstlisting 环境读取代码并应用全局 \lstset 设置
        \begin{lstlisting}[language=Python,
            caption={快速排序实现},
            label={lst:quicksort}]
def quicksort(arr):
    """快速排序算法"""
    if len(arr) <= 1:
        return arr
    pivot = arr[len(arr) // 2]
    left  = [x for x in arr if x < pivot]
    middle = [x for x in arr if x == pivot]
    right = [x for x in arr if x > pivot]
    return quicksort(left) + middle + quicksort(right)
        \end{lstlisting}
    \end{block}
\end{frame}

% ---- 6.2 Verilog 代码 ----
%  [fragile] 必需：因为 frame 中使用了 lstlisting 环境
\begin{frame}[fragile]{Verilog 代码展示}{IC 设计常用语言}
    %
    \begin{lstlisting}[language=Verilog,
        caption={D 触发器 Verilog 实现},
        label={lst:dff}]
module dff (
    input  wire clk,
    input  wire rst_n,
    input  wire d,
    output reg  q
);
    always @(posedge clk or negedge rst_n) begin
        if (!rst_n)
            q <= 1'b0;   // 异步复位
        else
            q <= d;       // 数据赋值
    end
endmodule
    \end{lstlisting}
\end{frame}

% ---- 6.3 行内代码与 verbatim ----
%  [fragile] 必需：因为 frame 中使用了 \verb 命令
\begin{frame}[fragile]{行内代码与 Verbatim}{小型代码片段}
    %
    \begin{itemize}
        \item 行内代码：\verb|\includegraphics[width=0.5\textwidth]{img.png}|
        \item 使用 \verb|\verb| 可以原样输出特殊字符
        \item \verb|\verb| 的定界符可以是任意相同字符，如 \verb|!$%^!|
              % 注意：上面的示例用 ! 而不是 | 作为定界符
        \item 在 Beamer 中，\verb|\verb| 不能用在 frame 的参数中，
              但可以安全地在 frame 体内部使用
    \end{itemize}

    \begin{warningblock}
        \verb|\verb| 命令不能用在 \verb|\section{}|、\verb|\caption{}|
        等移动参数中，也不能用在其他命令的参数内。
    \end{warningblock}
\end{frame}

% ============================================================
%  SECTION 7：TikZ 矢量绘图
% ============================================================
\section{TikZ 矢量绘图}

% ---- 7.1 基本图形 ----
%  [fragile] 必需：因为 frame 中使用了 \verb 命令
\begin{frame}[fragile]{TikZ 基本图形}{形状、颜色与线条}
    %
    \begin{columns}[T]
        \begin{column}{0.5\textwidth}
            \centering
            \begin{tikzpicture}[scale=0.9]
                % 矩形
                \fill[blue!20]  (0,2) rectangle (2,3);
                \draw[blue,thick] (0,2) rectangle (2,3);
                \node at (1,2.5) {矩形};

                % 圆形
                \fill[red!20]   (3.5,2.5) circle (0.7);
                \draw[red,thick] (3.5,2.5) circle (0.7);
                \node at (3.5,2.5) {圆形};

                % 椭圆
                \fill[green!20] (1,0.5) ellipse (1.5 and 0.6);
                \draw[green,thick] (1,0.5) ellipse (1.5 and 0.6);
                \node at (1,0.5) {椭圆};
            \end{tikzpicture}
        \end{column}
        \begin{column}{0.5\textwidth}
            \textbf{常用绘图元素}：
            \begin{itemize}
                \item \verb|\draw| —— 画线/边框
                \item \verb|\fill| —— 填充区域
                \item \verb|\node| —— 放置文本
                \item \verb|\coordinate| —— 定义坐标点
                \item \verb|rectangle| —— 矩形
                \item \verb|circle| —— 圆形
                \item \verb|ellipse| —— 椭圆
                \item \verb|arc| —— 圆弧
                \item \verb|--| —— 直线
                \item \verb|to [bend left]| —— 曲线
            \end{itemize}
        \end{column}
    \end{columns}
\end{frame}

% ---- 7.2 流程图 ----
\begin{frame}{TikZ 流程图}{节点、箭头与定位}
    %
    \centering
    \begin{tikzpicture}[
        node distance=1.5cm and 1.2cm,  % 节点间距（垂直 and 水平）
        box/.style={                    % 自定义样式 "box"
            rectangle,
            draw=darkblue,
            fill=blue!5,
            thick,
            minimum width=2.2cm,
            minimum height=0.8cm,
            rounded corners=3pt,
            align=center
        },
        arrow/.style={                  % 自定义样式 "arrow"
            ->,
            >=Stealth,
            thick,
            darkblue
        }
    ]
        % 定义节点（相对定位）
        \node[box] (A) {需求分析};
        \node[box, right=of A] (B) {方案设计};
        \node[box, below=of B] (C) {代码实现};
        \node[box, left=of C] (D) {测试验证};
        \node[box, below=of C] (E) {部署上线};

        % 画箭头
        \draw[arrow] (A) -- (B);
        \draw[arrow] (B) -- (C);
        \draw[arrow] (C) -- (D);
        \draw[arrow] (D) -- (A);
        \draw[arrow] (C) -- (E);
    \end{tikzpicture}
    %
    \vspace{0.2cm}
    \small 使用 \texttt{positioning} 库实现相对定位，\texttt{arrows.meta} 实现箭头样式。
\end{frame}

% ---- 7.3 简单数据图 ----
\begin{frame}{TikZ 数据示意图}{柱状图}
    %
    \centering
    \begin{tikzpicture}[scale=0.9]
        % Y 轴
        \draw[->] (0,0) -- (0,4.5) node[above] {性能指数};
        % X 轴
        \draw[->] (0,0) -- (6.5,0) node[right] {方案};

        % 柱状条
        \foreach \x/\y/\c/\l in {
            1/2.5/blue!40/A,
            2/3.8/red!40/B,
            3/3.0/green!40/C,
            4/4.0/orange!40/D
        }{
            \fill[\c] (\x*1.2-0.3, 0) rectangle (\x*1.2+0.3, \y);
            \node at (\x*1.2, -0.3) {\l};
            \node at (\x*1.2, \y+0.2) {\y};
        }
        % 虚线参考线
        \draw[dashed,gray] (0,3.5) -- (5.8,3.5) node[left,black] {3.5};
    \end{tikzpicture}
\end{frame}

% ============================================================
%  SECTION 8：动画与叠层（Overlays）
% ============================================================
\section{动画与叠层}

% ---- 8.1 pause 命令 ----
\begin{frame}{逐项显示：pause 命令}{最简单的动画方式}
    %
    %  \pause 是最简单的叠层命令，每遇到一个 \pause
    %  就把之后的内容放到下一页幻灯片中。
    %
    \begin{itemize}
        \item 第一项：这是第一页就显示的内容
        \pause
        \item 第二项：翻页后显示此项
        \pause
        \item 第三项：再翻页后显示此项
        \pause
        \item 第四项：最终显示此项
    \end{itemize}

    \pause
    \alert{每按一次空格/方向键，就显示下一个 \texttt{\textbackslash pause} 后的内容。}
\end{frame}

% ---- 8.2 叠层规范 ----
\begin{frame}{叠层规范}{精细控制显示时机}
    %
    %  叠层规范（Overlay Specification）使用尖括号 <...>，
    %  可以精确控制每条内容出现在第几张叠层中。
    %
    \begin{itemize}
        \item<1-> 这张幻灯片共有 5 个叠层
        \item<2-> <2-> 表示从第 2 层开始显示，之后一直可见
        \item<3-4> <3-4> 表示只在第 3、4 层可见
        \item<4-|alert@5> 第 4 层出现，第 5 层高亮（alert）
        \item<5-> 到这里全部呈现完毕
    \end{itemize}

    \uncover<3->{%
        %  \uncover<...> 控制整段内容的显示
        \begin{block}{动态 Block}
            这个 block 从第 3 层开始才显示。
        \end{block}
    }

    \only<5>{%
        %  \only<...> 的内容只在指定层出现，不占位
        \begin{center}
            \textcolor{darkred}{\textbf{—— 演示结束 ——}}
        \end{center}
    }
\end{frame}

% ---- 8.3 常用叠层修饰 ----
\begin{frame}{叠层的高级用法}{onslide / only / visible / invisible}
    %
    \begin{block}{各种叠层命令对比}
        \begin{description}
            \item[\texttt{\textbackslash pause}] 最简单的逐步显示
            \item[\texttt{\textbackslash onslide<2->}] 在指定叠层显示（占位不占位？需结合使用）
            \item[\texttt{\textbackslash only<2->}] 只在指定叠层显示（\important{不占空间}）
            \item[\texttt{\textbackslash uncover<2->}] 在指定叠层显示（\important{占空间}，留白）
            \item[\texttt{\textbackslash visible<2->}] 与 uncover 类似
            \item[\texttt{\textbackslash invisible<2->}] 在指定叠层隐藏
            \item[\texttt{\textbackslash alt<2>}] 不同叠层显示不同内容
        \end{description}
    \end{block}

    %  onslide 的实际演示
    \begin{columns}[T]
        \begin{column}{0.5\textwidth}
            \onslide<2->{%
                \textcolor{darkgreen}{从第 2 层开始显示此文本。}
            }
        \end{column}
        \begin{column}{0.5\textwidth}
            \only<3>{第 3 层：这段文字出现}
            \only<4>{第 4 层：这段文字替换了上面}
            \only<5>{第 5 层：再次替换}
        \end{column}
    \end{columns}
\end{frame}

% ---- 8.4 列表逐项显示 ----
\begin{frame}{列表的逐项显示}{itemize 与 enumerate 的叠层}
    %
    \begin{itemize}[<+->]
        % [<+->] 等价于给每个 \item 自动加上 <1->, <2->, ...
        \item 使用 \texttt{[<+->]} 选项可以让列表自动逐步显示
        \item 不需要在每个 \texttt{\textbackslash item} 后手动加 \texttt{\textbackslash pause}
        \item 这样代码更加简洁
        \item 但所有 item 会均匀占据空间（提前占位）
    \end{itemize}

    \vspace{0.3cm}
    对比（手动指定叠层顺序）：
    \begin{itemize}
        \item<5-> 此项在第 5 层才出现
        \item<6-> 此项在第 6 层才出现
        \item<7-> 此项在第 7 层才出现
    \end{itemize}
\end{frame}

% ============================================================
%  SECTION 9：引用、超链接与参考文献
% ============================================================
\section{引用与参考文献}

% ---- 9.1 超链接 ----
%  [fragile] 必需：因为 frame 中使用了 \verb 命令
\begin{frame}[fragile]{超链接}{内部与外部链接}
    %
    \begin{itemize}
        \item 外部链接：\href{https://www.latex-project.org}{\latex 官方网站}
        \item 外部 URL：\url{https://ctan.org/pkg/beamer}
        \item 内部跳转（到标签 \ref{sec:tips}）：
              见第 \ref{sec:tips} 节
        \item PDF 书签由 \verb|hyperref| + \verb|bookmark| 自动生成
    \end{itemize}

    \begin{tip}
        所有链接在 PDF 阅读器中均可点击跳转。
    \end{tip}
\end{frame}

% ---- 9.2 参考文献 ----
%  [fragile] 必需：因为 frame 中使用了 \verb 命令
\begin{frame}[fragile]{参考文献}{BibTeX / biblatex 引用}
    %
    在 Beamer 中引用参考文献同样方便：

    \begin{itemize}
        \item 使用 BibTeX：\verb|\bibliographystyle{ieeetr}| + \verb|\bibliography{refs}|
        \item 使用 biblatex：在导言区加载 \verb|biblatex| 宏包
        \item 引用命令：\verb|\cite{key}| 产生编号引用
        \item 在 frame 中用 \verb|\footcite{key}| 做脚注引用
    \end{itemize}

    \begin{block}{示例引用格式}
        \begin{itemize}
            \item Beamer 用户手册~\cite{beamer-manual}
            \item \latex 入门指南~\cite{latex-guide}
            \item TikZ \& PGF 手册~\cite{tikz-manual}
        \end{itemize}
    \end{block}
    % 注意：以下 \cite 引用的是虚拟 key，实际使用时需要对应的 .bib 文件
\end{frame}

% ============================================================
%  SECTION 10：实用技巧与最佳实践（带 label）
% ============================================================
\section{实用技巧与最佳实践}
\label{sec:tips}  % 用于前文的交叉引用

\begin{frame}{Beamer 最佳实践}{高效制作专业幻灯片的建议}
    %
    \begin{enumerate}[<+->]
        \item \textbf{内容为王}：每页幻灯片只传达 1--2 个核心观点
        \item \textbf{少即是多}：避免大段文字，多用列表和图表
        \item \textbf{统一风格}：全文使用一致的字体、颜色和布局
        \item \textbf{善用 Block}：用不同的 block 区分信息的性质
        \item \textbf{合理使用动画}：叠层有助于逐步呈现复杂概念，但不要过度
        \item \textbf{测试编译}：提交前用 \xelatex 编译两次确保交叉引用正确
        \item \textbf{备用方案}：同时导出 PDF 作为兜底（不依赖特定阅读器动画）
    \end{enumerate}
\end{frame}

\begin{frame}[fragile]{常见编译问题与解决}
    %  [fragile] 必需：enumerate 中有 \verb 命令
    \begin{warningblock}
        \textbf{\xelatex 编译常见问题}：
        \begin{enumerate}
            \item \textbf{中文不显示}：检查是否用了 \xelatex（不是 pdflatex）
            \item \textbf{字体找不到}：根据操作系统更换中文字体名
            \item \textbf{引用显示 ??}：需要编译两次
            \item \textbf{Overfull box}：内容超出页宽，考虑换行或缩小字号
            \item \textbf{lstlisting 出错}：fragile frame 需要加 \verb|[fragile]| 选项
        \end{enumerate}
    \end{warningblock}

    \begin{tip}
        结论：含有 \verb|\verb|、\verb|lstlisting| 等内容的帧，
        \important{必须在 \texttt{\textbackslash begin\{frame\}} 后加 \texttt{[fragile]} 选项}。
    \end{tip}
\end{frame}

% ============================================================
%  结束页
% ============================================================
\begin{frame}[plain]  % [plain] 去掉顶部和底部装饰
    \centering
    \vspace{2cm}

    {\Huge \textbf{谢谢！}}  % \Huge 是最大号字体命令

    \vspace{1cm}
    {\Large 问题与讨论}

    \vspace{1.5cm}
    {\normalsize
        GitHub Copilot \\
        \href{mailto:example@example.com}{example@example.com} \\
        \vspace{0.3cm}
        本模板源码可在以下路径找到：\\
        \texttt{content/template/beamer-template.tex}
    }

    \vspace{1cm}
    {\small 编译于 \today}
\end{frame}

% ============================================================
%  附录（Appendix）：额外的参考幻灯片
%  附录中的幻灯片不计入主进度条的总页数
% ============================================================
\appendix
%  \appendix 之后的所有 \section 都会自动标记为"附录"

\begin{frame}[allowframebreaks]{参考文献}
    %
    %  allowframebreaks 允许帧内容过长时自动分页
    %
    %  如果使用 BibTeX，可以这样：
    % \bibliographystyle{ieeetr}
    % \bibliography{references}
    %
    %  手动列出参考文献（模拟）：
    \begin{thebibliography}{99}
        \bibitem{beamer-manual} Till Tantau et al.,
            \emph{The Beamer Class: User Guide for version 3.x},
            \url{https://ctan.org/pkg/beamer}.

        \bibitem{latex-guide} Leslie Lamport,
            \emph{\latex: A Document Preparation System},
            Addison-Wesley, 1994.

        \bibitem{tikz-manual} Till Tantau,
            \emph{The TikZ and PGF Packages},
            \url{https://ctan.org/pkg/pgf}.
    \end{thebibliography}
\end{frame}

\begin{frame}{主题与配色速查表}
    %
    \begin{columns}[T]
        \begin{column}{0.33\textwidth}
            \textbf{整体主题}：
            \begin{itemize}
                \footnotesize
                \item Madrid
                \item Berlin
                \item Copenhagen
                \item Singapore
                \item Warsaw
                \item AnnArbor
                \item Boadilla
            \end{itemize}
        \end{column}
        \begin{column}{0.33\textwidth}
            \textbf{配色主题}：
            \begin{itemize}
                \footnotesize
                \item whale
                \item dolphin
                \item seahorse
                \item orchid
                \item crane
                \item beaver
            \end{itemize}
        \end{column}
        \begin{column}{0.33\textwidth}
            \textbf{字体主题}：
            \begin{itemize}
                \footnotesize
                \item default
                \item professionalfonts
                \item serif
                \item structurebold
            \end{itemize}
        \end{column}
    \end{columns}
\end{frame}

% ============================================================
%  文档结束
% ============================================================
\end{document}
